% ===========================================================================
%  Capítulo 1 — Preliminares teóricos
% ===========================================================================
\chapter{Preliminares te\'oricos}
\label{cap:teoria}

En este cap\'{\i}tulo se presentan los conceptos y herramientas que fundamentan
el trabajo.  La secci\'on~\ref{sec:hoja-calculo} define los t\'erminos clave
del dominio de las hojas de c\'alculo.  La secci\'on~\ref{sec:planificacion}
sit\'ua ese dominio en el contexto acad\'emico.  Las secciones~\ref{sec:dsl}
y~\ref{sec:ast} introducen los lenguajes de dominio espec\'{\i}fico y los
\'arboles de sintaxis abstracta.  Finalmente, la secci\'on~\ref{sec:lisp}
describe las caracter\'{\i}sticas de Common Lisp que lo convierten en la
plataforma de implementaci\'on elegida.

% ---------------------------------------------------------------------------
\section{Hojas de c\'alculo}
\label{sec:hoja-calculo}
% ---------------------------------------------------------------------------

Una \textbf{hoja de c\'alculo} (\textit{spreadsheet}) es una aplicaci\'on
para organizar, calcular y visualizar datos en una cuadr\'{\i}cula
bidimensional de celdas.  Cada celda puede contener un valor literal
---un n\'umero, una fecha o una cadena de texto--- o una
\textbf{f\'ormula}: una expresi\'on cuyo resultado se recalcula
autom\'aticamente cuando cambia cualquier celda de la que depende~\cite{microsoft-excel}.

\begin{description}
  \item[Workbook (libro de trabajo).] Un archivo \xlsx{} que agrupa una
    o m\'as hojas de trabajo.  Es la unidad de intercambio del formato
    Office Open XML~\cite{ooxml}.

  \item[Hoja de trabajo.] Una cuadr\'{\i}cula de celdas dentro del workbook,
    identificada por un nombre (p.~ej., \texttt{Defensas}).

  \item[Celda.] La unidad m\'{\i}nima de almacenamiento, referenciada por
    una direcci\'on del tipo \texttt{B4} (columna B, fila 4) o
    \texttt{\$B\$4} (referencia absoluta).

  \item[F\'ormula.] Una expresi\'on que comienza con \texttt{=} y puede
    contener referencias a otras celdas, rangos y funciones predefinidas
    como \texttt{COUNTIF}, \texttt{COUNTA}, \texttt{SUM}, \texttt{IF} o
    \texttt{SUMPRODUCT}.

  \item[Rango.] Un bloque rectangular de celdas contiguas, expresado como
    \texttt{\$B\$4:\$F\$6} (columnas B--F, filas 4--6).

  \item[Formato condicional.] Un mecanismo que aplica estilos visuales
    ---color de relleno, color de texto, borde--- a las celdas de un rango
    cuando se cumple una condici\'on evaluada en tiempo de apertura o
    edici\'on del archivo.  Una \textbf{FormulaRule} es el tipo m\'as general
    de regla de formato condicional: la condici\'on es una f\'ormula
    arbitraria, lo que permite expresar condiciones que involucran otras
    hojas, otras celdas o funciones de agregaci\'on~\cite{openpyxl}.
\end{description}

La biblioteca \textbf{openpyxl} es la herramienta Python elegida para
producir archivos \xlsx{} en este trabajo.  Permite crear workbooks,
escribir valores y f\'ormulas en celdas, y definir \texttt{FormulaRule}
que quedan incrustadas en el archivo y se eval\'uan en Excel o LibreOffice
Calc sin necesidad de c\'odigo adicional~\cite{openpyxl}.

% ---------------------------------------------------------------------------
\section{Planificaci\'on tabular en el \'ambito acad\'emico}
\label{sec:planificacion}
% ---------------------------------------------------------------------------

La planificaci\'on tabular agrupa los problemas de organizaci\'on cuya
soluci\'on se representa de forma natural como una o m\'as tablas
relacionadas.  En el contexto acad\'emico, los ejemplos m\'as frecuentes son:

\begin{itemize}
  \item \textbf{Horarios de docencia.} Distribuci\'on de grupos, asignaturas,
    profesores y aulas en slots de d\'{\i}a y hora a lo largo de la semana.
    Cada columna calcula la ocupaci\'on de un aula a partir de las asignaciones
    de los distintos grupos.

  \item \textbf{Planificaci\'on de defensas de tesis.} Asignaci\'on de
    estudiantes, fechas, horarios y miembros de tribunal (tutor, oponente,
    presidente, vocal, secretario) a defensas individuales.  El requisito
    cr\'{\i}tico es detectar conflictos: un mismo profesor en dos defensas
    del mismo slot horario.

  \item \textbf{Asignaci\'on de docencia.} Relaci\'on entre profesores,
    asignaturas y grupos, con conteos de horas acumuladas y alertas cuando
    se supera un umbral.
\end{itemize}

Estos documentos comparten una estructura t\'{\i}pica~\cite{ldmag2025}:
\begin{enumerate}
  \item Dos o m\'as tablas relacionadas que comparten columnas de clave
    (p.~ej., nombre de profesor, slot horario).
  \item Columnas calculadas en una tabla que agregan datos de otra
    (conteos, sumas, b\'usquedas).
  \item Reglas de formato condicional que colorean filas cuando se detecta
    una condici\'on problem\'atica (conflicto de horario, desequilibrio
    de carga).
  \item Totales y etiquetas fuera de las tablas, cuya posici\'on depende del
    n\'umero de filas de datos, que cambia en cada ejecuci\'on.
\end{enumerate}

Estos cuatro patrones son los que el DSL debe capturar de forma declarativa.

% ---------------------------------------------------------------------------
\section{Lenguajes de dominio espec\'{\i}fico}
\label{sec:dsl}
% ---------------------------------------------------------------------------

Un \textbf{lenguaje de dominio espec\'{\i}fico} (\textit{Domain-Specific
Language}, DSL) es un lenguaje de programaci\'on o especificaci\'on dise\~nado
para expresar soluciones dentro de un dominio de problemas concreto, en
contraposici\'on a los lenguajes de prop\'osito general que aspiran a cubrir
cualquier problema~\cite{fowler2010}.

\begin{definition}[DSL, \cite{van-deursen2000}]
Un lenguaje de dominio espec\'{\i}fico es un lenguaje de programaci\'on o
especificaci\'on ejecutable que ofrece, mediante abstracciones apropiadas,
poder expresivo centrado en y normalmente restringido a un dominio de
problemas particular.
\end{definition}

Los DSL se clasifican en dos grandes tipos seg\'un su relaci\'on con el
lenguaje anfitri\'on~\cite{fowler2010}:

\begin{description}
  \item[DSL externo.] Tiene su propia sintaxis independiente, requiere un
    lexer y un parser propios, y suele estar embebido como datos en otro
    programa (p.~ej., SQL, HTML, \LaTeX).

  \item[DSL interno (embebido).] Se expresa directamente en el lenguaje
    anfitri\'on mediante una API o un sistema de macros que produce c\'odigo
    v\'alido del anfitri\'on~\cite{mernik2005}.  El programador usa el
    lenguaje anfitri\'on como plataforma; el DSL es una biblioteca con
    vocabulario de dominio.
\end{description}

El DSL desarrollado en este trabajo es \textbf{interno}: est\'a escrito en
Common Lisp, sus formas son expresiones Lisp v\'alidas y no requiere ning\'un
parser externo.  Los macros de Lisp operan como el compilador del DSL,
expandiendo las formas del dominio en constructores de nodos del AST.

La ventaja principal de los DSL internos en Lisp es que el lenguaje anfitri\'on
es homoic\'onico: el c\'odigo tiene la misma estructura que los datos
(listas anidadas), lo que hace que escribir transformaciones de c\'odigo
---es decir, macros--- sea tan natural como escribir funciones ordinarias~\cite{graham1993}.

% ---------------------------------------------------------------------------
\section{\'Arboles de sintaxis abstracta}
\label{sec:ast}
% ---------------------------------------------------------------------------

Un \textbf{\'arbol de sintaxis abstracta} (AST, del ingl\'es
\textit{Abstract Syntax Tree}) es la representaci\'on intermedia est\'andar
en los compiladores e int\'erpretes.  Representa la estructura sint\'actica
de un programa mediante un \'arbol cuyos nodos corresponden a los constructores
del lenguaje y cuyos hijos representan los subcomponentes de cada
constructor~\cite{aho2006}.

El AST es \emph{abstracto} en el sentido de que omite los detalles de la
sintaxis concreta ---par\'entesis, palabras clave, signos de puntuaci\'on---
y conserva \'unicamente la estructura l\'ogica del programa~\cite{scott2015}.

En el dise\~no del DSL presentado en este trabajo, el AST cumple el papel
de \textbf{representaci\'on intermedia independiente del backend}: el
programador construye el AST mediante las formas del DSL, y el backend
lo recorre para producir el c\'odigo Python de generaci\'on del \xlsx{}.
Gracias a esta separaci\'on, el mismo AST podr\'{\i}a en principio alimentar
backends alternativos (LibreOffice Calc, CSV, HTML) sin modificar la capa
del DSL.

Cada nodo del AST es una instancia de una clase CLOS definida con la macro
\dsl{defclass*}.  Los nodos se dividen en cuatro familias:
\begin{enumerate}
  \item Nodos de estructura del workbook (\dsl{xl-workbook}, \dsl{xl-sheet},
    \dsl{xl-region}, \dsl{xl-table}).
  \item Nodos de expresi\'on de celda (\dsl{xl-expr-column-ref},
    \dsl{xl-table-range}, \dsl{xl-expr-countif}, \dsl{xl-expr-if}, etc.).
  \item Nodos de cuantificaci\'on existencial (\dsl{xl-expr-exists},
    \dsl{xl-domain-rows}).
  \item Nodos de formato condicional y posicionamiento (\dsl{xl-style-rule},
    \dsl{xl-anchor}, \dsl{xl-nav}, \dsl{xl-sheet-fixed-expr}).
\end{enumerate}

Esta organizaci\'on se describe en detalle en el
cap\'{\i}tulo~\ref{cap:dsl}.

% ---------------------------------------------------------------------------
\section{El lenguaje de programaci\'on Common Lisp}
\label{sec:lisp}
% ---------------------------------------------------------------------------

Common Lisp es un lenguaje de programaci\'on de prop\'osito general creado
originalmente por John McCarthy en 1958~\cite{mccarthy1960} y estandarizado
como Common Lisp en 1994~\cite{steele1990}.  Su elecci\'on como plataforma
de implementaci\'on del DSL se justifica por cuatro caracter\'{\i}sticas
decisivas.

% ............................................................................
\subsection{Notaci\'on prefija y homoiconicidad}
\label{ssec:prefija}
% ............................................................................

En Common Lisp, todas las expresiones son listas anidadas con notaci\'on
prefija: el primer elemento es el operador y los siguientes son los
argumentos.  Una llamada a funci\'on, la definici\'on de una clase, una
iteraci\'on y una expresi\'on del DSL tienen exactamente la misma forma
sint\'actica.

Esta propiedad, llamada \textbf{homoiconicidad}, hace que el c\'odigo Lisp
sea indistinguible de los datos Lisp~\cite{graham1993}.  El compilador puede
inspeccionar y transformar el c\'odigo del mismo modo que inspecciona y
transforma listas de datos.  Los DSL internos se benefician directamente
de esto: las formas del DSL son listas ordinarias que los macros transforman
en c\'odigo v\'alido.

% ............................................................................
\subsection{Macros}
\label{ssec:macros}
% ............................................................................

Un \textbf{macro} de Lisp es una funci\'on que opera sobre el c\'odigo fuente
del programa ---representado como listas--- y devuelve c\'odigo nuevo que
sustituye a la forma original antes de la compilaci\'on o evaluaci\'on~\cite{seibel2005}.

La diferencia con una funci\'on ordinaria es temporal: una funci\'on se
eval\'ua en tiempo de ejecuci\'on; un macro se expande en tiempo de
compilaci\'on.  Esta distinci\'on permite a los macros introducir nueva
sintaxis y nuevas abstracciones sin coste en tiempo de ejecuci\'on.

En este trabajo, cada constructo del DSL ---\dsl{def-table},
\dsl{conditional-rendering}, \dsl{exists}, \dsl{nav}--- se implementa como
un macro que construye uno o m\'as nodos del AST.  El programador escribe en
vocabulario del dominio; el macro transforma esa escritura en llamadas a los
constructores CLOS.

% ............................................................................
\subsection{El Sistema de Objetos de Common Lisp (CLOS)}
\label{ssec:clos}
% ............................................................................

El \textbf{Sistema de Objetos de Common Lisp} (CLOS, del ingl\'es
\textit{Common Lisp Object System}) es el sistema de orientaci\'on a objetos
integrado en el est\'andar del lenguaje~\cite{clos1988}.  Sus caracter\'{\i}sticas
principales son:

\begin{description}
  \item[Herencia m\'ultiple.] Una clase puede heredar de varias superclases
    simult\'aneamente.  Los conflictos de herencia se resuelven mediante el
    \textit{Class Precedence List} (CPL), un orden lineal determinista~\cite{keene1989}.

  \item[Despacho m\'ultiple.] Los m\'etodos no pertenecen a una clase
    espec\'{\i}fica.  Una funci\'on gen\'erica puede despachar sobre los tipos
    de varios argumentos simult\'aneamente~\cite{kiczales1991}.  Esto significa
    que el c\'omportamiento de \dsl{(compile-excel-formula nodo ...)} depende
    del tipo de \dsl{nodo}, no de un m\'etodo ligado a una clase particular.

  \item[Funciones gen\'ericas.] Las funciones polim\'orficas se declaran con
    \dsl{defgeneric} y se especializan con \dsl{defmethod}.  A\~nadir soporte
    para un nuevo tipo de nodo AST s\'olo requiere a\~nadir un nuevo
    \dsl{defmethod}; el resto del sistema no se modifica.
\end{description}

En este trabajo, CLOS cumple dos funciones.  Primero, las clases CLOS
representan los nodos del AST: cada tipo de nodo es una clase con slots
que almacenan sus atributos.  Segundo, la funci\'on gen\'erica
\dsl{compile-excel-formula} usa el despacho m\'ultiple para compilar
cada tipo de nodo a su cadena de f\'ormula Excel correspondiente, sin
ninguna estructura \dsl{cond} o \dsl{case} centralizada.

% ............................................................................
\subsection{La macro \texttt{defclass*}}
\label{ssec:defclass-star}
% ............................................................................

Para reducir el c\'odigo repetitivo de definir clases CLOS y sus
constructores, el sistema define la macro \dsl{defclass*}, que genera
simult\'aneamente:

\begin{itemize}
  \item Una clase CLOS con nombre \dsl{clase-xl-\textit{nombre}}, donde
    cada slot tiene un accessor y un \texttt{initarg} del mismo nombre.
  \item Una funci\'on constructora \dsl{xl-\textit{nombre}} con
    argumentos por clave (\dsl{\&key}), que crea una instancia de la clase
    con los valores proporcionados.
\end{itemize}

\begin{lstlisting}[language=CommonLisp,
  caption={Implementaci\'on de defclass*.}]
(defmacro defclass* (class-name parents slots
                      &key (ctr-args-type '&key))
  `(progn
     (defclass ,(symb 'clase- class-name) ,parents
       ,(loop for s in slots
              collecting (expansion-del-slot s)))
     ,(make-constructor class-name slots
                        :argument-type ctr-args-type)))
\end{lstlisting}

Por ejemplo, la definici\'on del nodo \dsl{xl-expr-if}:

\begin{lstlisting}[language=CommonLisp,
  caption={Ejemplo de uso de defclass*.}]
(defclass* xl-expr-if () (test then else))
;; genera:
;;   (defclass clase-xl-expr-if () (test :accessor test :initarg :test)
;;                                  (then :accessor then :initarg :then)
;;                                  (else :accessor else :initarg :else))
;;   (defun xl-expr-if (&key test then else)
;;     (make-instance 'clase-xl-expr-if :test test :then then :else else))
\end{lstlisting}

Este patr\'on elimina el c\'odigo repetitivo de definici\'on de nodos y
garantiza que cada clase tiene una funci\'on constructora consistente con
su interfaz de inicializaci\'on.
